\documentclass[11pt,a4paper]{article}

\usepackage[T1]{fontenc}
\usepackage[utf8]{inputenc}
\usepackage[brazil]{babel}
\usepackage{amsmath,amssymb}
\usepackage[margin=2.5cm]{geometry}
\usepackage{enumitem}
\usepackage{tikz}
\usetikzlibrary{arrows.meta,positioning}

\setlist[enumerate]{leftmargin=*}

\newcommand{\Enc}{\mathsf{Enc}}
\newcommand{\Dec}{\mathsf{Dec}}
\newcommand{\Gen}{\mathsf{Gen}}
\newcommand{\Mac}{\mathsf{Mac}}
\newcommand{\Vrfy}{\mathsf{Vrfy}}
\newcommand{\negl}{\mathsf{negl}}
\newcommand{\Ruim}{\mathsf{Ruim}}
\newcommand{\Repete}{\mathsf{Repete}}
\newcommand{\ilegivel}[1]{\textit{[ilegível: #1]}}
\newcommand{\defeq}{\stackrel{\text{\tiny def}}{=}}
\newcommand{\PrivKcpaI}[3]{\mathsf{PrivK}^{\mathrm{cpa}\text{-}1}_{#1,#2}(#3)}
\newcommand{\PrivKcpa}[3]{\mathsf{PrivK}^{\mathrm{cpa}}_{#1,#2}(#3)}
\newcommand{\PrivKbit}[3]{\mathsf{PrivK}^{\mathrm{cpa}\text{-}1}_{#1,#2}(#3)}
\newcommand{\PrivKccaI}[3]{\mathsf{PrivK}^{\mathrm{cca}\text{-}1}_{#1,#2}(#3)}

\title{Lista 2 \\ EDUARDO KALUF - GRR20241770 \\ PEDRO TAKEO SHIMA - GRR20240627}
\author{}
\date{}

\begin{document}
    \maketitle

%======================================================================
    \section*{Questão 1 --- ``Seja $F\colon \{0,1\}^{n} \times \{0,1\}^{n} \to \{0,1\}^{n}$
        uma função pseudoaleatória (chaveada). Neste exercício, vamos mostrar que nenhum
        adversário probabilístico polinomial $\mathcal{A}$ consegue encontrar dois valores
        $x$ e $y$ tal que $F_k(x) = F_k(y)$ com probabilidade não-negligenciável.''}
%======================================================================

    \subsection*{a) ``Seja $\mathcal{A}$ um algoritmo de tempo polinomial $q(n)$ tal que
        $\mathcal{A}^{\mathcal{O}(\cdot)}(1^{n}) = 1$ se e somente se $\mathcal{A}$ encontra
        uma colisão na função $\mathcal{O}\colon \{0,1\}^{n} \to \{0,1\}^{n}$.
        Argumente que, se $f\colon \{0,1\}^{n} \to \{0,1\}^{n}$ é uma
        função aleatória, então:
        \[
            \Pr\!\left[\mathcal{A}^{f(\cdot)}(1^{n}) = 1\right] \le \frac{q(n)^{2}}{2^{n+1}}.\text{''}
        \]
    }

% Resolução

    \subsection*{b) ``Usando o fato de que $F$ é pseudoaleatória, conclua que existe uma
    função negligenciável $\negl$ tal que
        \[
            \Pr\!\left[\mathcal{A}^{F_k(\cdot)}(1^{n}) = 1\right] \le \negl(n).\text{''}
        \]}

% Resolução

    \clearpage

%======================================================================
    \section*{Questão 2 --- ``Prove que o modo de operação CTR não é CCA-seguro.''}
%======================================================================

% Resolução

    \clearpage

%======================================================================
    \section*{Questão 3 --- ``(Exercício 4.1 do livro, adaptado) Considere um MAC
        $\Pi = (\Gen, \Mac, \Vrfy)$ tal que, para uma chave $k \in \{0,1\}^{n}$, $\Mac_k$
        sempre tem como saída uma tag de tamanho $c \cdot \log_2 n$, onde $c$ é uma
        constante. Prove que $\Pi$ não pode ser um MAC seguro.''}
%======================================================================

% Resolução

    \clearpage

%======================================================================
    \section*{Questão 4 --- ``Considere uma variação de segurança CCA (que chamamos de
    segurança CCA-1) em que o adversário $\mathcal{A}$ só pode fazer consultas ao
    oráculo de decifragem antes de enviar um par de mensagens para o desafiante.
    Chamaremos este experimento de $\PrivKccaI{\mathcal{A}}{\Pi}{n}$. Neste exercício,
        vamos provar que a Construção 3.30, além de CPA-segura, ainda é CCA-1-segura.''}
%======================================================================

    \subsection*{a)}

    \begin{quote}
        ``Seja $\tilde{\Pi}$ uma versão da cifra da Construção 3.30 onde se usa uma função
        aleatória no lugar da função pseudoaleatória $F$. Fixe um adversário probabilístico
        polinomial $\mathcal{A}$ e considere o experimento
        $\PrivKccaI{\mathcal{A}}{\tilde{\Pi}}{n}$. Defina $r^{*} \in \{0,1\}^{n}$ como a
        string aleatória utilizada pelo desafiante para cifrar a mensagem de desafio $m_b$.
        Por fim, defina o evento $\Ruim$ como a união de dois eventos $\Repete$ e
        $\Repete'$, definidos a seguir:

        \begin{itemize}
            \item $\Repete$ é o evento de que, no experimento
            $\PrivKccaI{\mathcal{A}}{\tilde{\Pi}}{n}$, o desafiante usa $r^{*}$ como string
            aleatória para cifrar alguma mensagem consultada por $\mathcal{A}$, antes ou
            depois de $\mathcal{A}$ receber o texto cifrado de desafio do desafiante. Esse é
            o mesmo evento considerado na demonstração da aula 10.
            \item $\Repete'$ é o evento de que, no experimento
            $\PrivKccaI{\mathcal{A}}{\tilde{\Pi}}{n}$, o adversário envia para o desafiante
            alguma consulta de decifragem do tipo $\langle r^{*}, \cdot \rangle$ antes de
            receber o texto cifrado de desafio do desafiante.
        \end{itemize}

        Prove que $\Pr[\Ruim] \le \dfrac{q(n)}{2^{n-1}}$ e que
        $\Pr\!\left[\PrivKccaI{\mathcal{A}}{\tilde{\Pi}}{n} = 1 \mid \overline{\Ruim}\right]
        = \dfrac{1}{2}$, onde $q(n)$ é um limite para o número de consultas de
        cifragem/decifragem que $\mathcal{A}$ realiza. Conclua então que
        \[
            \Pr\!\left[\PrivKccaI{\mathcal{A}}{\tilde{\Pi}}{n} = 1\right]
            \le \frac{1}{2} + \frac{q(n)}{2^{n-1}}.\text{''}
        \]
    \end{quote}

% Resolução

    \subsection*{b)}

    \begin{quote}
        ``Usando o fato de que $F$ é uma função pseudoaleatória, prove que existe uma função
        negligenciável $\negl$ tal que:
        \[
            \left|\,\Pr\!\left[\PrivKccaI{\mathcal{A}}{\Pi}{n} = 1\right]
                - \Pr\!\left[\PrivKccaI{\mathcal{A}}{\tilde{\Pi}}{n} = 1\right]\right|
            \le \negl(n).
        \]
        Conclua então que a Construção 3.30 é CCA-1-segura.''
    \end{quote}

% Resolução

    \clearpage

%======================================================================
    \section*{Questão 5 --- ``(Exercício 5.3 do livro, adaptado) Seja $(\Gen, H)$ uma
    função hash resistente à colisões, e defina $(\Gen, \hat{H})$ como
        $\hat{H}^{s}(x) \defeq H^{s}(H^{s}(x))$. A função $(\Gen, \hat{H})$ é também
        resistente à colisões? Justifique sua resposta.''}
%======================================================================

% Resolução

\end{document}